\documentclass[../../main.tex]{subfiles}
% \documentclass[a4paper,10pt]{ltjsarticle}
\begin{document}
\rhead{京都大学後期$1993$年$4$番}
\setcounter{equation}{0}
\begin{oframed}
  \subfile{problem}
\end{oframed}
\setlength{\columnseprule}{0.4pt}
\begin{multicols}{2}
  {\bf [解答]}

  \begin{figure}[H]
    \centering
    \begin{tikzpicture}[y=0.3cm, >=stealth, every node/.style={font=\small}]
      % 左：a=3/2 の例．1<x<beta (beta≒4.94) で a^x < ax となり条件を満たさない．
      \begin{scope}[x=0.55cm]
        \pgfmathsetmacro{\bt}{4.939}
        \fill[pattern=north east lines, pattern color=red!75!black]
        plot[domain=1:\bt, samples=60] (\x, {1.5*\x})
        -- plot[domain=\bt:1, samples=60] (\x, {1.5^\x}) -- cycle;
        \draw[->, thick, gray!80!black] (-0.4, 0) -- (6.0, 0) node[right] {$x$};
        \draw[->, thick, gray!80!black] (0, -0.8) -- (0, 10.2) node[above] {$y$};
        \draw[line width=1.4pt, blue!85!black, domain=-0.4:5.6, samples=80]
        plot (\x, {1.5^\x}) node[above left] {$y=a^x$};
        \draw[line width=1.4pt, teal!80!black, domain=-0.3:5.6]
        plot (\x, {1.5*\x}) node[right] {$y=ax$};
        \draw[densely dotted, thick, gray!70!black] (1, 1.5) -- (1, 0) node[below] {$1$};
        \draw[densely dotted, thick, gray!70!black] (\bt, {1.5*\bt}) -- (\bt, 0) node[below] {$\beta$};
        \node[circle, fill=red!85!black, inner sep=1.4pt] at (1, 1.5) {};
        \node[circle, fill=red!85!black, inner sep=1.4pt] at (\bt, {1.5*\bt}) {};
        \node[below left=1pt, font=\footnotesize] at (0, 0) {$O$};
        \node[below=16pt, font=\footnotesize] at (2.8, 0) {(i) $a=\frac{3}{2}$ のとき};
      \end{scope}
      % 右：a=e のとき x=1 で接し，全ての x>0 で a^x >= ax．
      \begin{scope}[x=1.0cm, xshift=4.6cm]
        \draw[->, thick, gray!80!black] (-0.3, 0) -- (2.2, 0) node[right] {$x$};
        \draw[->, thick, gray!80!black] (0, -0.8) -- (0, 10.2) node[above] {$y$};
        \draw[line width=1.4pt, blue!85!black, domain=-0.3:2.2, samples=80]
        plot (\x, {exp(\x)}) node[left] {$y=a^x$};
        \draw[line width=1.4pt, teal!80!black, domain=-0.15:2.2]
        plot (\x, {exp(1)*\x}) node[below right=-1pt] {$y=ax$};
        \draw[densely dotted, thick, gray!70!black] (1, {exp(1)}) -- (1, 0) node[below] {$1$};
        \draw[densely dotted, thick, gray!70!black] (1, {exp(1)}) -- (0, {exp(1)}) node[left] {$e$};
        \node[circle, fill=red!85!black, inner sep=1.4pt] at (1, {exp(1)}) {};
        \node[below left=1pt, font=\footnotesize] at (0, 0) {$O$};
        \node[below=16pt, font=\footnotesize] at (1, 0) {(ii) $a=e$ のとき};
      \end{scope}
    \end{tikzpicture}
    \caption{$y=a^x$と$y=ax$のグラフ（斜線部では$a^x<ax$）}
    \label{1993-4:fig:1}
  \end{figure}

  $f(x) = a^x - ax$ とおく．$x>0$ に対し $f(x)\ge 0$ となるための条件を求めれば良い．
  ただし，題意より
  \begin{align}
    a > 0
  \end{align}
  である．
  $0 < a \le 1$ とすると $f(x) \to -\infty$ ($x \to \infty$) から不適．
  従って
  \begin{align}
    1 < a \label{eq:1}
  \end{align}
  として考える．

  $f'(x)$の一階微分は
  \begin{align}
    f'(x) = a \left[ a^{x-1} \log a - 1 \right]
  \end{align}
  だから，$f'(x)=0$の解は
  \begin{align}
    a^{\alpha-1} = \dfrac{1}{\log a} \label{eq:2}
  \end{align}
  なる$\alpha$である．$a>1$よりこの両辺は正なので自然対数をとって，
  \begin{align}
    (\alpha-1) \log a & = - \log (\log a)                                     \\
    \therefore
    \alpha            & = 1 - \dfrac{\log(\log a)}{\log a} \label{eq:3}
  \end{align}
  $f(x)$の増減表は\cref{1993-4:table:1}となる．

  \begin{table}[H]
    \centering
    \caption{$f(x)$の増減表}
    \label{1993-4:table:1}
    \begin{tabular}{|c||c|c|c|c|}
      \hline
      $x$  & $0$ & $\cdots$   & $\alpha$ &            \\
      \hline
      $f'$ &     & $-$        & $0$      & $+$        \\
      \hline
      $f$  &     & $\searrow$ &          & $\nearrow$ \\
      \hline
    \end{tabular}
  \end{table}

  従って，題意の条件を満たすには $f(\alpha) \ge 0$ ならば良い．
  \begin{align}
    f(\alpha)
    & = a\left[a^{\alpha-1}-\alpha\right] \\
    & = a\left[\dfrac{1}{\log a} -\alpha\right] \quad (\because\text{\cref{eq:2}})
  \end{align}
  だから，
  \begin{align}
    & f(\alpha) \ge 0              \\
    \iff
    & \alpha \le \dfrac{1}{\log a }
  \end{align}
  \cref{eq:3}を代入して
  \begin{align}
    & 1 - \dfrac{\log(\log a)}{\log a} \le \dfrac{1}{\log a } \\
    & \log\dfrac{a}{\log a } \le 1
  \end{align}
  なる条件を得る．

  $\log x$ は $x$ について単調増加なので，$\log e = 1$に注意して$\log$を外して
  \begin{align}
    & \dfrac{a}{\log a} \le  e \\
    \therefore
    & e\log a -a \ge 0
  \end{align}
  となる．以下この条件を満たすような$a$を求めるため，
  \begin{align}
    g(a) = e\log a -a
  \end{align}
  とおくと，$g(a)\ge 0$となる$a$を求めれば良い．
  一階微分は
  \begin{align}
    g'(a) = \dfrac{e}{a} -1
  \end{align}
  だから，$g(x)$の増減表は\cref{1993-4:table:2}となる．

  \begin{table}[H]
    \centering
    \caption{$g(x)$の増減表}
    \label{1993-4:table:2}
    \begin{tabular}{|c||c|c|c|c|}
      \hline
      $a$  & $1$ &            & $e$ &            \\
      \hline
      $g'$ &     & $+$        &     & $-$        \\
      \hline
      $g$  &     & $\nearrow$ & $0$ & $\searrow$ \\
      \hline
    \end{tabular}
  \end{table}
  従って，$g(a) \ge 0$を満たす $a$ は $a=e$のみであり，これは\cref{eq:1}を満たすから，
  求める$a$は
  \begin{align}
    a = e
  \end{align}
  である． $\cdots$(答)

  {\bf [解説]}
  解答ではかなり素直に関数を微分して条件を満たす$a$を求めに行ったが，よりショートカットも可能ある．
  題意の条件
  \begin{align}
    a^x \ge ax
  \end{align}
  が任意の$x\ge 0$で満たすような$a$を探せば良い．$a>1$よりこの両辺は正なので$a$を底とする対数をとって
  \begin{align}
    x-1 \ge \log_a x
  \end{align}
  となるから，
  \begin{align}
    h(x) = x-1 - \log_a x
  \end{align}
  の増減表を書いて$h(x)\ge 0$となる条件を求めにいけば良い．

  \newpage
\end{multicols}
\end{document}
